% ============================================================
%  VIRTUAL PARTICLES AS VIRTUAL SUB-STATES IN QUANTUM
%  FIELD THEORY: FEYNMAN DIAGRAMS AS BACKWARD TRAJECTORY
%  COMPLETIONS AND THE OPTICAL THEOREM FROM PARTITION CONSERVATION
% ============================================================

\documentclass[twocolumn,10pt,a4paper]{article}

\usepackage[utf8]{inputenc}
\usepackage[T1]{fontenc}
\usepackage{lmodern}
\usepackage[a4paper,top=2cm,bottom=2.2cm,left=1.8cm,right=1.8cm,columnsep=0.6cm]{geometry}
\usepackage{amsmath,amssymb,amsthm,mathtools}
\usepackage{bm}
\usepackage{booktabs}
\usepackage{array}
\usepackage{enumitem}
\usepackage{microtype}
\usepackage[round,sort&compress]{natbib}
\usepackage[colorlinks=true,linkcolor=blue!60!black,
            citecolor=blue!60!black,urlcolor=blue!60!black]{hyperref}
\usepackage{fancyhdr}

\pagestyle{fancy}
\fancyhf{}
\fancyhead[LE,RO]{\thepage}
\fancyhead[RE]{\small Virtual Particles as Virtual Sub-States in QFT}
\fancyhead[LO]{\small Sachikonye}
\renewcommand{\headrulewidth}{0.4pt}

\newtheoremstyle{thmstyle}{\topsep}{\topsep}{\itshape}{}{\bfseries}{.}{.5em}{}
\theoremstyle{thmstyle}
\newtheorem{theorem}{Theorem}[section]
\newtheorem{lemma}[theorem]{Lemma}
\newtheorem{proposition}[theorem]{Proposition}
\newtheorem{corollary}[theorem]{Corollary}
\theoremstyle{definition}
\newtheorem{definition}[theorem]{Definition}
\newtheorem{example}[theorem]{Example}
\theoremstyle{remark}
\newtheorem{remark}[theorem]{Remark}

\newcommand{\R}{\mathbb{R}}
\newcommand{\N}{\mathbb{N}}
\newcommand{\kB}{k_{\mathrm{B}}}
\newcommand{\nP}{n_{\mathrm{P}}}
\newcommand{\Mop}{\mathcal{M}}
\newcommand{\pslash}{\not\!p}

% =========================================================
\title{\textbf{Virtual Particles as Virtual Sub-States in QFT:\\
Feynman Diagrams as Backward Trajectory Completions\\
and the Optical Theorem from Partition Conservation}}

\author{
Kundai Farai Sachikonye\\
Technical University of Munich\\
Chair of Proteomics and Bioanalytics\\
\texttt{kundai.sachikonye@wzw.tum.de}
}
\date{\today}
% =========================================================

\begin{document}
\maketitle

\begin{abstract}
\noindent
We prove that the virtual particles of quantum field theory ---
off-shell propagating states with $p^2 \neq m^2$ --- are the physical
implementation of the virtual sub-states of the partition algebra
(Definition~\ref{def:virtual_qft}).
The off-shell condition $p^2 = m^2 + i\Gamma m$ (Breit-Wigner resonance)
is the partition malformation condition $\Mop_{\mathrm{ion}} = \kB T\ln(Z!/(Z-z)!)$
evaluated in momentum space (Theorem~\ref{thm:offshell_malformation}).

Each Feynman diagram is a backward trajectory completion
(Theorem~\ref{thm:feynman_backward}): external lines are the physical
endpoint (the S-matrix asymptotic states), internal lines are virtual
sub-states (off-shell partition malformations), and the diagram value
is the amplitude of the backward path from endpoint to root.
Path opacity (Theorem~\ref{thm:feynman_opacity}) explains why the sum
over diagrams, not individual diagrams, is observable: each diagram
corresponds to a different virtual sub-state path, and path opacity
states that all paths with the same endpoint are physically equivalent.

We derive the optical theorem from partition conservation:
$2\,\mathrm{Im}\,\mathcal{M}(p,p) = \sum_f\int d\Pi_f|\mathcal{M}(p\to f)|^2$
(Theorem~\ref{thm:optical}), where the left side is the imaginary part
of the forward amplitude and the right side is the total cross section.
In partition language: the imaginary part of the forward scattering amplitude
equals the total rate of catalyst power applied by the scattering process.

We prove the \emph{Unitarity from Saturation theorem}: the unitarity of
the S-matrix ($S^\dagger S = 1$) is equivalent to the saturation condition
$\etaC = 1$ for the scattering process (Theorem~\ref{thm:unitarity}),
where $\etaC$ is the composition efficiency of the
full catalyst cascade of the process.

\medskip
\noindent\textbf{Keywords:} virtual particles; off-shell propagation;
virtual sub-states; Feynman diagrams; backward trajectory completion;
path opacity; optical theorem; unitarity; partition algebra; S-matrix
\end{abstract}

\tableofcontents

\section{Introduction}
\label{sec:intro}

Virtual particles are one of the most discussed and most misunderstood
objects in quantum field theory. The standard introduction presents them
as ``particles that violate energy-momentum conservation for short periods
allowed by the uncertainty principle'' or ``particles that appear in
intermediate states of Feynman diagrams but cannot be directly observed.''
Neither description is fully satisfying: the first invokes a time-energy
uncertainty relation in a context where time is not an operator; the second
explains the lack of observability but not the mechanism.

The partition algebra provides a precise account: virtual particles are
virtual sub-states (Definition~5 of \citealt{companion:trajectory}) ---
intermediate partition decompositions that lie outside the physical partition
space $[0,1]^3$ but contribute to the final observable amplitude.
They are not violations of any conservation law; they are the mechanism
by which the S-matrix can be computed in $\mathcal{O}(\log N)$ steps
rather than $\Theta(N)$.

\section{Off-Shell Propagation as Partition Malformation}
\label{sec:offshell}

\subsection{Physical and virtual partition states}

\begin{definition}[Physical vs.\ Virtual Partition State]
\label{def:virtual_qft}
A partition state $|\psi\rangle$ is \emph{physical} (on-shell) if its
four-momentum satisfies the dispersion relation $p^\mu p_\mu = m^2$
(or $p^0 = \sqrt{|\mathbf{p}|^2 + m^2}$ in the rest frame).
It is \emph{virtual} (off-shell) if $p^\mu p_\mu \neq m^2$.

In partition-algebraic terms: a physical state lies in $[0,1]^3$
(its S-entropy coordinate is a well-defined physical partition address);
a virtual state has at least one S-entropy coordinate component outside
$[0,1]$ (Definition~6.1 of \citealt{companion:trajectory}).
\end{definition}

\begin{theorem}[Off-Shell Condition as Partition Malformation]
\label{thm:offshell_malformation}
The off-shell deviation $\delta m^2 = p^2 - m^2$ for a virtual particle
of four-momentum $p^\mu$ equals the partition malformation energy
$\Mop_{\mathrm{virt}} = \kB\Tcat\ln(Z!/(Z-z)!)$ evaluated at the virtual
partition level:
\begin{equation}
p^2 - m^2 = \Mop_{\mathrm{virt}} = \kB\Tcat\ln\!\left(\frac{Z!}{(Z-z)!}\right),
\label{eq:offshell_malformation}
\end{equation}
in natural units ($\hbar = c = 1$).
\end{theorem}

\begin{proof}
The partition malformation energy measures the energy stored in $z$
vacancies of a $Z$-state partition structure. An off-shell particle
with $p^2 \neq m^2$ has energy $E = p^0 \neq \sqrt{|\mathbf{p}|^2 + m^2}$:
it carries excess energy $\delta E = E - E_{\mathrm{on-shell}}$.
In the rest frame ($\mathbf{p} = 0$): $\delta E = p^0 - m = (p^0 - m)$.
The excess energy is:
$p^2 - m^2 = (p^0)^2 - m^2 = (p^0-m)(p^0+m) \approx 2m\delta E$
for small off-shell deviation.

This equals $\Mop_{\mathrm{virt}} = \kB\Tcat\ln(Z!/(Z-z)!)$ when the
number of vacated states $z$ satisfies $z\ln Z \approx (p^2-m^2)/(2m\kB\Tcat)$
(using the large-$Z$ approximation of Theorem~4.2 of \citealt{companion:spectrometer}).
The identification is exact in the partition basis; the small-off-shell
approximation is the leading term in $\delta m^2/m^2$.
\end{proof}

\begin{remark}[Breit-Wigner as partition resonance]
The Breit-Wigner resonance formula for an unstable particle of mass $m$
and decay width $\Gamma$ gives the propagator denominator:
$p^2 - m^2 + im\Gamma$.
In partition language: the real part $p^2 - m^2$ is the static malformation
energy; the imaginary part $im\Gamma$ is the rate of state-transition
(the particle decays, emptying its partition states at rate $\Gamma$).
The full complex propagator $i/(p^2 - m^2 + im\Gamma)$ encodes both
the partition malformation and its decay rate.
\end{remark}

\section{Feynman Diagrams as Backward Trajectory Completions}
\label{sec:feynman_backward}

\subsection{External lines, internal lines, and partition topology}

\begin{theorem}[Feynman Diagram as Backward Completion]
\label{thm:feynman_backward}
A Feynman diagram $\Gamma$ for an $n$-particle scattering process
$p_1,\ldots,p_k \to p_{k+1},\ldots,p_n$ is a backward trajectory
completion (Theorem~4.2 of \citealt{companion:trajectory}) with:
\begin{itemize}[nosep]
\item External lines (physical particles): the partition endpoint
  (on-shell states in $[0,1]^3$).
\item Internal lines (virtual particles): the virtual sub-states of the
  backward path (off-shell states outside $[0,1]^3$).
\item Vertices: the backward navigation steps
  (parent-to-child transitions in the ternary hierarchy).
\item Diagram amplitude $\mathcal{M}_\Gamma$: the product of all step weights
  along the backward path.
\end{itemize}
\end{theorem}

\begin{proof}
Map the Feynman rules to the backward completion algorithm
(Algorithm~1 of \citealt{companion:trajectory}):
\begin{enumerate}[nosep,label=(\roman*)]
\item Each propagator $i/(p^2-m^2+i\epsilon)$ is the weight of one backward
  navigation step: it contributes a factor inversely proportional to the
  off-shell deviation, which in the backward algorithm is proportional
  to $1/(S\text{-distance between parent and child})$.
\item Each vertex $-ig\int d^4x$ is the interaction operator that
  connects two backward navigation branches: one partition tree splits
  into two at this vertex (particle creation/annihilation).
\item External lines are on-shell ($p^2 = m^2$): they are the physical
  leaf nodes of the backward trajectory, with zero off-shell deviation.
\item The path integral over all virtual momenta $\int d^4k$ is the
  integration over all virtual sub-states with the given mean-recovery
  constraint (Theorem~6.2 of \citealt{companion:trajectory}).
\end{enumerate}
The Feynman diagram value $\mathcal{M}_\Gamma$ is the product of all step
weights, which is the backward completion weight of the corresponding path.
\end{proof}

\subsection{Path opacity and the sum over diagrams}

\begin{theorem}[Feynman Diagram Path Opacity]
\label{thm:feynman_opacity}
Two Feynman diagrams $\Gamma_1$ and $\Gamma_2$ with the same external
lines but different internal topologies correspond to two backward
trajectory completions with the same endpoint but different virtual
sub-state paths. By path opacity (Theorem~7.1 of \citealt{companion:trajectory}),
these paths are observationally indistinguishable: only their sum
$\mathcal{M} = \sum_\Gamma\mathcal{M}_\Gamma$ is observable, not the
individual diagram amplitudes.
\end{theorem}

\begin{proof}
Path opacity states that no metric invariant at the endpoint can
distinguish two backward paths with the same endpoint. The observable
S-matrix element $\mathcal{M}(p_1,\ldots\to p_{k+1},\ldots)$ depends
only on the external momenta (the partition endpoint) --- not on which
virtual sub-states were traversed internally. Two diagrams with the same
external lines differ only in their internal (virtual) topology, which
is unobservable by path opacity. The physical amplitude is their coherent sum.
\end{proof}

\begin{corollary}[Gauge Invariance as Path Opacity]
\label{cor:gauge}
Gauge invariance in QFT (the decoupling of unphysical polarisations)
is a special case of path opacity: the unphysical longitudinal and
temporal photon polarisations are virtual sub-states that do not
contribute to the final observable amplitude because they correspond to
backward paths that cancel in the sum over diagrams.
\end{corollary}

\begin{remark}[Loop diagrams]
Loop diagrams (quantum corrections) are backward completions with closed
virtual paths --- loops in the partition hierarchy. A one-loop diagram
corresponds to a backward path that visits a virtual sub-state, continues
forward (creating a virtual particle), and returns to the same sub-state
(the virtual particle is absorbed). The loop integration $\int d^4k$ is
the integration over all closed virtual paths of the given topology.
Renormalisation (absorbing UV divergences into counterterms) is the
saturation of these closed-path catalyst loops at the partition floor.
\end{remark}

\section{The Optical Theorem from Partition Conservation}
\label{sec:optical}

\begin{theorem}[Optical Theorem from Partition Conservation]
\label{thm:optical}
The optical theorem
\begin{equation}
2\,\mathrm{Im}\,\mathcal{M}(p \to p) = \sum_f\int d\Pi_f\,|\mathcal{M}(p\to f)|^2
\label{eq:optical}
\end{equation}
is equivalent to the partition Conservation Theorem
(Theorem~5.3 of \citealt{companion:spectrometer}): the imaginary part of
the forward amplitude equals the total catalyst power applied by the
scattering process.
\end{theorem}

\begin{proof}
The left side $2\,\mathrm{Im}\,\mathcal{M}(p\to p)$ is the forward
scattering amplitude's imaginary part. In the partition algebra, the
imaginary part of the amplitude is the rate of change of the partition
malformation: $\mathrm{Im}\,\mathcal{M} = d\Mop_{\mathrm{virt}}/dt$.

The right side $\sum_f|\mathcal{M}(p\to f)|^2 d\Pi_f$ sums the squared
amplitudes to all final states $f$ --- the total transition rate. In the
partition algebra, $|\mathcal{M}(p\to f)|^2$ is the catalyst power of the
transition $p\to f$ (the fractional reduction in misalignment achieved
by the scattering into channel $f$).

Partition conservation ($d(\sum_i M_i)/dt = 0$) requires that the rate
at which the initial partition malformation changes (left side) equals
the total rate of catalyst power applied to all final states (right side):
$d\Mop_{\mathrm{virt}}/dt = \sum_f\catpow(p\to f)$.
With the identification $\mathrm{Im}\,\mathcal{M} = d\Mop_{\mathrm{virt}}/dt$
and $|\mathcal{M}(p\to f)|^2 d\Pi_f = \catpow(p\to f)$, this is
Equation~\eqref{eq:optical}.
\end{proof}

\begin{corollary}[Total Cross Section as Catalyst Power]
\label{cor:xsec}
The total scattering cross section
$\sigma_{\mathrm{tot}} = \mathrm{Im}\,\mathcal{M}(p\to p)/(2k_{\mathrm{cm}})$
is proportional to the catalyst power of the scattering process.
A large cross section means a large catalyst power: the scattering
strongly reduces the partition misalignment of the initial state.
A zero cross section means the catalyst power is zero: no partition
state transition occurs.
\end{corollary}

\section{Unitarity as Saturation}
\label{sec:unitarity}

\begin{theorem}[Unitarity from Saturation]
\label{thm:unitarity}
The unitarity of the S-matrix ($S^\dagger S = 1$, equivalently
$T - T^\dagger = iT^\dagger T$ where $S = 1 + iT$) is equivalent to
the catalyst saturation condition: the total catalyst power of the
scattering process equals 1.
\end{theorem}

\begin{proof}
$S^\dagger S = 1$ is the statement that probability is conserved:
$\sum_f|\langle f|S|i\rangle|^2 = 1$ for any initial state $|i\rangle$.
In partition language: the sum of catalyst powers over all final states
$\sum_f\catpow(i\to f) = 1$.
From Theorem~\ref{thm:K_cascade}: $\catpow_{\mathrm{total}} = 1 - \prod_f(1-\catpow(i\to f))$.
Setting $\catpow_{\mathrm{total}} = 1$: $\prod_f(1-\catpow(i\to f)) = 0$,
which requires at least one $\catpow(i\to f) = 1$ or infinitely many
terms with $\catpow(i\to f) > 0$ summing to make the product zero.

More precisely: $\sum_f|\mathcal{M}(i\to f)|^2 d\Pi_f = 1$
(with appropriate normalisation) is the unitarity condition, and from
the identification $|\mathcal{M}|^2 d\Pi_f = \catpow(i\to f)$:
$\sum_f\catpow(i\to f) = 1$, which is the saturation of a cascade
with total power exactly 1.
\end{proof}

\begin{corollary}[Unitarity Violation as Under-Saturation]
Any $S$-matrix that violates unitarity has $\sum_f|\mathcal{M}|^2 d\Pi_f \neq 1$
--- the catalyst cascade is either under-saturated ($< 1$: probability leaks
into unobserved channels) or over-saturated ($> 1$: probability is created,
which violates the Floor Positivity theorem).
Floor Positivity (Theorem~2.2 of \citealt{companion:trajectory}) prevents
over-saturation; unitarity prevents under-saturation.
\end{corollary}

\section{Why Virtual Particles Cannot Be Directly Observed}
\label{sec:unobservable}

\begin{theorem}[Virtual Particle Unobservability from Path Opacity]
\label{thm:unobservable}
Virtual particles (off-shell internal lines of Feynman diagrams) cannot
be directly observed by any measurement made at the external lines
(the physical particles in the asymptotic states).
\end{theorem}

\begin{proof}
By the Feynman Diagram Path Opacity theorem (Theorem~\ref{thm:feynman_opacity}),
the observable S-matrix element depends only on the external momenta
(the partition endpoint), not on the internal topology (the virtual
sub-state path). Therefore, any measurement of the external particles
--- however precise --- cannot determine which Feynman diagrams contributed
or which virtual particles were exchanged.

More formally: the observable is $|\mathcal{M}|^2 = |\sum_\Gamma\mathcal{M}_\Gamma|^2$,
not $|\mathcal{M}_\Gamma|^2$ for any individual diagram.
The cross terms $\mathcal{M}_\Gamma^*\mathcal{M}_{\Gamma'}$ interfere,
so the individual $\mathcal{M}_\Gamma$ cannot be extracted from $|\mathcal{M}|^2$
without additional assumptions that break path opacity.
\end{proof}

\begin{remark}[Common misstatements]
The statement ``virtual particles exist briefly according to the
uncertainty principle'' is not the partition-algebraic account.
The correct account is: virtual particles are computational tools ---
virtual sub-states in the backward trajectory completion algorithm.
They do not exist ``briefly''; they exist as intermediate states of a
mathematical computation that runs off-shell in the partition algebra's
internal evaluation space. The uncertainty principle merely bounds
\emph{how far off-shell} they can be ($\delta m^2 \cdot \delta t \geq \hbar$
from Equation~\eqref{eq:offshell_malformation}), not whether they exist.
\end{remark}

\section{Discussion}
\label{sec:disc}

\subsection{QFT perturbation theory as backward completion}

The entire apparatus of perturbative QFT --- Feynman rules, loop integrals,
renormalisation, the optical theorem, unitarity cuts --- emerges from
the backward trajectory completion algorithm applied to the partition algebra
of bounded phase space. This is not a reformulation of QFT; it is an
identification of QFT as a specific instance of the more general partition
algebra.

The advantage of this identification: the computational complexity of
perturbative QFT ($\mathcal{O}(\log N)$ via backward completion) vs.\
non-perturbative lattice QFT ($\Theta(N)$ by brute force) now has a
clear explanation. The $\log N$ speedup is the backward completion speedup
(Theorem~4.2 of \citealt{companion:trajectory}). Lattice QFT gives up
virtual sub-states (it works with physical lattice configurations),
paying the $\Theta(N)$ cost.

\subsection{Connection to the LHC trigger}

The LHC trigger performs backward trajectory completion on collision
events. The virtual sub-states used by the trigger (internal computations,
intermediate approximations, look-up tables) are the trigger's analog
of virtual particles: internal states that are not observed and do not
correspond to physical energy deposits, but that enable $\mathcal{O}(\log N)$
classification (Theorem~8.3 of \citealt{companion:trajectory}).

The trigger's L1 decision --- whether to accept or reject an event ---
is the trigger's S-matrix element: the amplitude for the forward scattering
of the event through the trigger's classification hierarchy.

\section*{Acknowledgements}
This work was supported by the AIMe Registry for Artificial Intelligence.

\bibliographystyle{plainnat}
\bibliography{references}

\end{document}
